% ============================================================
% CUMCM 国赛 A 题机理类模板 — v7.2
% 编译:xelatex main.tex(跑两遍)
% 依赖:ctexart, amsmath, amssymb, graphicx, booktabs, hyperref
% 金标准:六段子结构 / 三段式公式 / 四重检验 / 算法对比 / 创新量化
% ============================================================
\documentclass[12pt,a4paper]{ctexart}

\usepackage{amsmath,amssymb,amsfonts}
\usepackage{graphicx}
\usepackage{booktabs}
\usepackage{geometry}
\usepackage{float}
\usepackage{caption}
\usepackage{subcaption}
\usepackage{listings}
\usepackage{hyperref}
\usepackage{enumitem}
\usepackage{xcolor}
\usepackage{indentfirst}
\usepackage{titlesec}
\usepackage{fancyhdr}
\usepackage{setspace}

\geometry{top=2.5cm, bottom=2.5cm, left=2.5cm, right=2.5cm, headheight=14pt}
% ============ 排版规范（对齐国赛优秀论文 B159） ============
% 行距 1.5 倍，首行缩进 2 字符
\linespread{1.5}\selectfont
\setlength{\parindent}{2em}

% 章节标题居中，二级标题左对齐
\titleformat{\section}{\centering\large\bfseries\heiti}{\thesection}{0.8em}{}
\titleformat{\subsection}{\normalsize\bfseries\heiti}{\thesubsection}{0.8em}{}
\titleformat{\subsubsection}{\normalsize\bfseries\heiti}{\thesubsubsection}{0.6em}{}

% 页眉空、页脚页码居中
\pagestyle{fancy}
\fancyhf{}
\fancyfoot[C]{\small\thepage}
\renewcommand{\headrulewidth}{0pt}


% 中文字体设置(评委机器通用)
\setCJKmainfont{SimSun}[BoldFont=SimHei, ItalicFont=KaiTi]
\setCJKsansfont{SimHei}
\setCJKmonofont{FangSong}

% 图/表题自动编号前缀（ctexart 默认不设置，缺失会导致图题无「图」字）
\renewcommand{\figurename}{图}
\renewcommand{\tablename}{表}


% 自定义算子(三段式公式规范)
\DeclareMathOperator{\erf}{erf}
\DeclareMathOperator{\sgn}{sgn}

% 超链接样式
\hypersetup{colorlinks=true, linkcolor=black, citecolor=black, urlcolor=blue}

% 代码块样式
\lstset{
  basicstyle=\ttfamily\small,
  breaklines=true,
  frame=single,
  numbers=left,
  numberstyle=\tiny,
  keywordstyle=\color{blue},
  commentstyle=\color{green!50!black},
}

\title{\textbf{[研究对象]的[核心问题]建模与求解}}
\author{队伍编号:\textbf{XXXX}}
\date{}

\begin{document}

\maketitle

% ============================================================
% 摘要三段式(半页内)
% ============================================================
\begin{abstract}
\noindent
\textbf{第一段(1-2 行)}:行业背景极度精简 + 研究对象 + 三层框架概括。

\textbf{第二段(核心)}:各问依次:模型+算法+精准量化结果,数值前置。例如:针对问题一,建立[模型],采用[算法]求解,得到[关键数值];针对问题二,...

\textbf{第三段(2-3 行)}:3 项创新(创新点1:精度提升 X\%;创新点2:速度提升 Y 倍;创新点3:稳定性提升 Z\%)+ 四重检验指标(R²/MAE/MC-CV)+ 跨行业迁移。

\vspace{0.5em}
\noindent\textbf{关键词}:[对象] [模型] [算法] [方法] [指标]
\end{abstract}

\newpage

% ============================================================
% 目录(编译两遍才能正确显示)
% ============================================================
\newpage

% ============================================================
% 一、问题重述
% ============================================================
\section{问题重述}

\subsection{问题背景}
[精简背景,1-2 段,聚焦研究对象]

\subsection{问题提出}
\begin{itemize}
  \item 问题一:[一句话描述]
  \item 问题二:[一句话描述]
  \item 问题三:[一句话描述]
\end{itemize}

% ============================================================
% 二、问题分析(附图1:物理过程示意图)
% ============================================================
\section{问题分析}

针对问题一,本文从[物理过程]角度出发,[分析思路]。

\begin{figure}[H]
\centering
\includegraphics[width=0.85\textwidth]{../figures/pdf/fig1_process.pdf}
\caption{图1 [物理过程]的[趋势]([关键数值])}
\label{fig:process}
\end{figure}

% ============================================================
% 三、模型假设(含物理简化假设)
% ============================================================
\section{模型假设}

\begin{enumerate}
  \item 假设1:[内容]
  \item 假设2:[内容]
  \item 假设3:[内容]
\end{enumerate}

% ============================================================
% 四、符号说明(物理量三线表)
% ============================================================
\section{符号说明}

\begin{table}[H]
\centering
\caption{符号说明表}
\label{tab:symbols}
\begin{tabular}{lllll}
\toprule
符号 & 物理量 & 数值/范围 & 单位 & 来源 \\
\midrule
$I_{\text{DNI}}$ & 直射辐照度 & 200-1000 & W/m² & 气象数据 \\
$\alpha$ & 太阳高度角 & 0-90 & ° & 计算 \\
$\theta$ & 太阳方位角 & 0-360 & ° & 计算 \\
$\eta_{\cos}$ & 余弦效率 & 0-1 & - & 计算 \\
\bottomrule
\end{tabular}
\end{table}

% ============================================================
% 五、模型建立与求解(每问六段子结构)
% ============================================================
\section{模型建立与求解}

% ---------- 问题一 ----------
\subsection{问题一}

% 5.1.1 参数配置与数据预处理
\subsubsection{参数配置与数据预处理}
[数据来源、字段说明、异常处理、参数取值]

% 5.1.2 基础经典模型(三段式公式)
\subsubsection{基础经典模型}

太阳光线经定日镜反射后汇聚于吸热塔顶端吸热口,反射过程遵循菲涅尔反射定律,需建立光线追踪模型。
\begin{equation}
\mathbf{r}_{\text{out}} = \mathbf{r}_{\text{in}} - 2(\mathbf{r}_{\text{in}} \cdot \mathbf{n})\mathbf{n}
\label{eq:reflection}
\end{equation}
其中 $\mathbf{r}_{\text{in}}$ 为入射光线单位向量,$\mathbf{r}_{\text{out}}$ 为反射光线单位向量,$\mathbf{n}$ 为镜面法向量。式~\eqref{eq:reflection} 表明反射光线关于法向量对称,需进一步计算吸热口截获效率。

% 5.1.3 传统模型缺陷剖析
\subsubsection{传统模型缺陷定量剖析}

传统二维投影法存在以下缺陷:
\begin{enumerate}
  \item 缺陷1:[内容],仿真数据显示误差达 $X\%$
  \item 缺陷2:[内容],在 [场景] 下失效
  \item 缺陷3:[内容],计算精度仅 $Y\%$
\end{enumerate}

% 5.1.4 创新改进模型
\subsubsection{自主创新改进模型}

为克服上述缺陷,本文建立基于三维光线追踪的改进模型,决策变量为 $\mathbf{x} = (r_i, \theta_i, h)$。
\begin{align}
\eta_{\cos}(t) &= \cos\theta_{\text{inc}}(t) \label{eq:cos_eff}\\
P_{\text{out}}(t) &= \eta_{\cos}(t) \cdot \eta_{\text{atm}}(t) \cdot I_{\text{DNI}}(t) \cdot A_i \label{eq:power}
\end{align}
其中 $\theta_{\text{inc}}$ 为入射角,$\eta_{\text{atm}}$ 为大气透射率。式~\eqref{eq:cos_eff} 给出余弦效率,式~\eqref{eq:power} 给出输出功率。相比式~\eqref{eq:reflection} 的二维方法,本模型精度提升 $12.3\%$。

% 5.1.5 算法设计与求解
\subsubsection{算法设计与求解}

采用 Runge-Kutta 四阶方法求解常微分方程组:
\begin{equation}
\frac{d\mathbf{x}}{dt} = \mathbf{f}(\mathbf{x}, t), \quad \mathbf{x}(0) = \mathbf{x}_0
\label{eq:ode}
\end{equation}
采用式~\eqref{eq:ode} 的四阶 Runge-Kutta 方法求解,步长 $h = 0.01$ s,精度 $10^{-6}$。

% 5.1.6 数值结果与可视化
\subsubsection{数值结果与可视化}

\begin{figure}[H]
\centering
\includegraphics[width=0.85\textwidth]{../figures/pdf/fig2_3d_surface.pdf}
\caption{图2 [对象]的[趋势](峰值 $Z_{\max} = 0.85$)}
\label{fig:3d_surface}
\end{figure}

\begin{table}[H]
\centering
\caption{问题一求解结果对比}
\label{tab:result1}
\begin{tabular}{lccc}
\toprule
指标 & 传统方法 & 改进方法 & 提升 \\
\midrule
精度 (\%) & 87.5 & 99.8 & +12.3 \\
计算时间 (s) & 45.2 & 0.53 & 85× \\
\bottomrule
\end{tabular}
\end{table}

% ---------- 问题二 ----------
\subsection{问题二}
% 同六段子结构,此处省略

% ---------- 问题三 ----------
\subsection{问题三}
% 同六段子结构,此处省略

% ============================================================
% 六、模型检验(四重检验,缺一降档)
% ============================================================
\section{模型检验}

\subsection{拟合精度检验}
$R^2 = 0.987$, $\text{MAE} = 0.23$, $\text{RMSE} = 0.31$, $\text{MAPE} = 2.3\%$。残差 Shapiro-Wilk 检验 $p = 0.87$,符合正态分布。

\begin{figure}[H]
\centering
\includegraphics[width=0.85\textwidth]{../figures/pdf/figN_residual.pdf}
\caption{图N 残差直方图与 Q-Q 图(残差符合正态分布)}
\label{fig:residual}
\end{figure}

\subsection{灵敏度分析}

\begin{table}[H]
\centering
\caption{参数灵敏度分析(弹性系数分级)}
\label{tab:sensitivity}
\begin{tabular}{lccccc}
\toprule
参数 & 基准值 & ±10\% 影响 & ±20\% 影响 & 弹性系数 $\varepsilon$ & 分级 \\
\midrule
镜宽 $w$ & 6 m & ±1.2\% & ±2.8\% & 0.14 & 低 \\
塔高 $h$ & 120 m & ±3.5\% & ±7.8\% & 0.39 & 中 \\
镜数 $N$ & 200 & ±5.2\% & ±11.3\% & 0.57 & 高 \\
\bottomrule
\end{tabular}
\end{table}

\subsection{蒙特卡洛验证}
$N_{\text{sim}} = 200$ 次,均值 $\mu = 53.8$ MW,标准差 $\sigma = 1.2$ MW,$\text{CV} = 2.2\%$,$95\%\text{CI} = [51.4, 56.2]$ MW。

\subsection{假设误差量化}

\begin{table}[H]
\centering
\caption{误差来源分解}
\label{tab:error}
\begin{tabular}{lccc}
\toprule
误差来源 & 估计值 & 占比 & 说明 \\
\midrule
数值离散化 & 0.032\% & 8\% & RK4 方法误差 \\
参数不确定性 & 0.25\% & 62\% & 输入参数测量误差 \\
模型简化 & 0.12\% & 30\% & 忽略高阶项 \\
\bottomrule
\end{tabular}
\end{table}

% ============================================================
% 七、模型优缺点与改进
% ============================================================
\section{模型优缺点与改进}

\textbf{优点}:
\begin{enumerate}
  \item 创新点1:建立三维光线追踪模型,相比传统二维投影法精度提升 $12.3\%$
  \item 创新点2:推导高效余弦效率解析表达式,计算速度提升 $85$ 倍
  \item 创新点3:提出分段自适应步长求解策略,数值稳定性提升 $40\%$
\end{enumerate}

\textbf{缺点}:对 [场景] 的适应性有限,未来可引入 [方法] 改进。

% ============================================================
% 八、参考文献
% ============================================================
\section{参考文献}

\begin{thebibliography}{99}
\bibitem{ref1} 作者. 标题[J]. 期刊, 年份, 卷(期): 页码.
\bibitem{ref2} Author A, Author B. Title[J]. Journal, Year, Vol(Issue): Pages.
% 注:近5年文献 ≥40%,外文 ≥30%,总数 ≥10
\end{thebibliography}

% ============================================================
% 附录(三模块)
% ============================================================
\appendix

\section{附录1 数据说明}
[数据来源、字段说明、result 样表 5 行 8 列]

\section{附录2 代码清单}
% ============ 国一铁律：附录必须含完整代码，缺代码取消评奖资格 ============
% 完整代码用 \lstinputlisting 从 code/ 目录读入，禁止只贴片段/伪代码。
% 每个脚本附输入/输出/功能说明，确保可复现（全部固定 SEED=42）。
\begin{table}[H]
\centering
\caption{代码模块说明（完整代码见附件，固定 SEED=42 可复现）}
\label{tab:code_modules}
\begin{tabular}{llll}
\toprule
脚本 & 功能 & 输入 & 输出 \\
\midrule
\texttt{code/00\_data\_prep.py} & 数据预处理与质量校验 & 附件1.xlsx & \texttt{data/clean/*.csv} \\
\texttt{code/01\_q1\_analysis.py} & 问题一求解 & 清洗后数据 & \texttt{results/q1\_*.json} \\
\texttt{code/02\_q2\_classify.py} & 问题二求解 & 关键词表 & \texttt{results/q2\_*.json} \\
\texttt{code/03\_q3\_optimize.py} & 问题三求解 & 问题二结果 & \texttt{results/q3\_*.json} \\
\texttt{code/04\_q4\_robust.py} & 问题四求解 & 问题三模型 & \texttt{results/q4\_*.json} \\
\texttt{code/05\_figures.py} & 图表生成 & \texttt{results/*.json} & \texttt{figures/png}, \texttt{figures/pdf} \\
\bottomrule
\end{tabular}
\end{table}

% 逐文件引入完整代码（示例）：
% \lstinputlisting[language=Python, caption=code/01_q1_analysis.py]{../code/01_q1_analysis.py}


\section{附录3 公式推导}
[核心公式完整推导,正文只放结论]

\end{document}
